Cette partie vise à déterminer le sens d'évolution d'un système chimique.

On mélange le même volume $V_{0}$ d'une solution aqueuse d'acide éthanoïque
$\ce{CH3CO2H_{(aq)}}$ et d'une solution aqueuse du benzoate de sodium
$\ce{C6H5CO^-_2_{(aq)} + Na^+_{(aq)}}$. Les deux solutions ont la même
concentration molaire $C_{0}$.
\begin{donnees}
\begin{itemize}
    \item $\mathrm{K_{A_1}}=
          \mathrm{K_{A}}(\ce{CH3CO2H_{(aq)}}\ttslash{}\ce{CH3CO^-_2_{(aq)}})=
          \num{1.8e-5}$~; 
    \item $\mathrm{K_{A_2}}=
          \mathrm{K_{A}}(\ce{C6H5CO2H_{(aq)}}\ttslash{}\ce{C6H5CO^-_2_{(aq)}})=
          \num{6.3e-5}$.
\end{itemize}
\end{donnees}

\begin{enumerate}
    \item Écrire l'équation chimique de la réaction qui se produit entre l'acide
          éthanoïque et l'ion benzoate.
    \item Montrer que l'expression de la constante d'équilibre $\mathrm{K}$
          associée à l'équation de cette réaction s'écrit
          $\mathrm{K}=\frac{\mathrm{K_{A_1}}}{\mathrm{K_{A_2}}}$, puis calculer
          sa valeur.
    \item La valeur du quotient de réaction du système chimique à l'état initial
          est $Q_{r,i}=1$.

          Dans quel sens évolue le système chimique~? Justifier.
\end{enumerate}

